\section{Shlosman's random staircases}

\begin{definition}[The discrete Gaussian staircase model, Georgii (6.16)--(6.20)]
    \label{def:rs-6-16}
    \uses{def:pot-nearest-neighbour, def:pot-gibbs-spec}
    \lean{Potential.discreteGaussian, MeasureTheory.GibbsMeasure.Shlosman.dgPotential, MeasureTheory.GibbsMeasure.Shlosman.isSigmaFiniteLambdaAdmissible_dgPotential, MeasureTheory.GibbsMeasure.Shlosman.dgSpecification, MeasureTheory.GibbsMeasure.Shlosman.map_dgPotential_eq, MeasureTheory.GibbsMeasure.Shlosman.staircase, MeasureTheory.GibbsMeasure.Shlosman.isGroundState_staircase}
    \leanok

    On $S = \mathbb Z^2$ with $E = \mathbb Z$ and counting measure, the potential with
    $H_\Lambda(\zeta) = \tfrac12\sum_{\langle i,j\rangle}(\zeta_i - \zeta_j)^2$ is
    $\lambda$-admissible for $\beta > 0$, invariant under the five symmetries (6.17)(i)--(v), and
    the staircases $\omega^z_i = z i_1 + i_2$, translated by any $n$, are its ground states.
\end{definition}

\begin{theorem}[Georgii (6.25)]
    \label{thm:rs-6-25}
    \uses{def:rs-6-16, thm:ising-betac-bounds}
    \lean{MeasureTheory.GibbsMeasure.Shlosman.exists_circuit_contour_dg, MeasureTheory.GibbsMeasure.Shlosman.stepDown, MeasureTheory.GibbsMeasure.Shlosman.card_bdIdx_le_hamiltonian_sub_stepDown, MeasureTheory.GibbsMeasure.Shlosman.dgSpecification_abs_excess_le}
    \leanok

    For $\beta > 0$, every $z$, every box $\Lambda_N$, every site $a$ and every $k$,
    $\gamma_{\Lambda_N}(\abs{\sigma_a - \omega^z_a} \geq k \mid \omega^z) \leq 2 r'(\beta/2)^k$,
    where $r'$ is the Peierls series of the sharp contour estimate.
\end{theorem}
\begin{proof}
    \leanok
    A configuration deviating at $a$ carries a circuit contour around $a$; the step-down map
    $t_c$ lowers the spins inside it, gains a factor $e^{-\beta}$ per contour bond by (6.24), and
    is injective, so the excess event is bounded by the anchored-circuit series; iterate in $k$
    and use the reflection symmetry to cover both signs.
\end{proof}

\begin{theorem}[Georgii (6.21)]
    \label{thm:rs-6-21}
    \uses{thm:rs-6-25, prop:cp-49}
    \lean{MeasureTheory.GibbsMeasure.Shlosman.staircasePhase, MeasureTheory.GibbsMeasure.Shlosman.staircasePhase_mem_GP, MeasureTheory.GibbsMeasure.Shlosman.half_lt_staircasePhase, MeasureTheory.GibbsMeasure.Shlosman.staircasePhase_agree_ge, MeasureTheory.GibbsMeasure.Shlosman.staircasePhase_ne, MeasureTheory.GibbsMeasure.Shlosman.infinite_GP_dgSpecification_of_log_twelve, MeasureTheory.GibbsMeasure.Shlosman.map_spinTranslation_staircasePhase_ne, MeasureTheory.GibbsMeasure.Shlosman.map_spinReflection_staircasePhase_ne, MeasureTheory.GibbsMeasure.Shlosman.map_shift_staircasePhase_ne, MeasureTheory.GibbsMeasure.Shlosman.map_latticeRot_staircasePhase_ne, MeasureTheory.GibbsMeasure.Shlosman.map_latticeReflFst_staircasePhase_ne, MeasureTheory.GibbsMeasure.locallyEquicontinuous_of_uniformlyTight}
    \leanok

    For $\beta \geq \log 12$ and every slope $z$ there is $\mu_z^\beta \in \mathcal G(\beta\Phi)$
    with $\mu_z^\beta(\sigma_a = \omega^z_a) > \tfrac12$ at every site and
    $\mu_z^\beta(\sigma_i = \omega^z_i \text{ on } \Lambda) \geq 1 - 2\abs{\Lambda}r'(\beta/2)$.
    Distinct slopes give distinct measures, and $\mu_z^\beta$ differs from its images under the
    spin translation, the spin reflection, the shifts and the lattice rotations and reflections;
    in particular $\mathcal G(\beta\Phi)$ is infinite and all five symmetries (6.17) are broken.
\end{theorem}
\begin{proof}
    \leanok
    Theorem \ref{thm:rs-6-25} makes the one-spin marginals uniformly tight, hence the finite-volume
    kernels locally equicontinuous, so a cluster point $\mu_z^\beta$ exists by Proposition
    \ref{prop:cp-49}; the same estimate transfers to it and pins the spin at each site to
    $\omega^z$ with probability above $\tfrac12$, which separates the phases.
\end{proof}
